\chapter{LA GRANDE CASA}\label{la-grande-casa}

\hfill\break

Avevo dodici anni e i gemelli tredici la notte in cui le stelle
scomparvero dal cielo.

Era ottobre, un paio di settimane prima di Halloween, e in occasione di
un ricevimento solo per adulti a noi tre era stato ordinato di rimanere
nel seminterrato di casa Lawton - la Grande Casa, la chiamavamo.

Essere confinati nel seminterrato non era una punizione. Non per Diane e
Jason, che per loro scelta passavano là quasi tutto il tempo; di certo
non lo era per me. Loro padre aveva stabilito un confine ben preciso tra
la zona degli adulti e quella dei bambini, ma nel nostro spazio c'erano
una piattaforma di gioco di ultima generazione, un'ampia videoteca e
persino un tavolo da biliardo\ldots{} e nessun controllo da parte degli
adulti, tranne quello della solita addetta al catering, una certa
signora Truall, che scendeva giù all'incirca ogni ora per evitarsi
l'incombenza di servire le tartine e darci aggiornamenti sulla festa.
(Un uomo della Hewlett-Packard aveva fatto una figuraccia con la moglie
di un commentatore del \emph{Post}. C'era un senatore ubriaco nello
studio.)

«Ci manca solo il silenzio,» disse Jason (l'impianto stereo del piano
superiore suonava musica da ballo che pulsava attraverso il soffitto
come il battito di un orco), «e la vista del cielo.»

Il silenzio e la vista: Jase, come suo solito, aveva deciso di volerli
entrambi.

Diane e Jason erano nati a qualche minuto di distanza ma era chiaro che
fossero fratelli dizigotici e non monozigotici; nessuno li chiamava
gemelli, tranne la madre. Jason diceva che erano il frutto di
«spermatozoi dipolari che avevano penetrato ovuli di carica opposta.»
Diane, il cui quoziente d'intelligenza era notevole quasi quanto quello
di Jason ma che superava il fratello in padronanza linguistica,
paragonava lei e Jase a «prigionieri diversi evasi dalla stessa cella.»

Avevo soggezione di entrambi.

Jason, a tredici anni, non era solo un mostro d'intelligenza; era anche
in perfetta forma fisica - non particolarmente muscoloso ma forte, e
spesso il migliore nelle gare di corsa. Già allora era alto quasi un
metro e ottanta, magro, la faccia goffa compensata da un sorriso
sbilanciato ma genuino. A quel tempo aveva i capelli biondi e ispidi.

Diane era tredici centimetri più bassa, rotondetta solo se paragonata al
fratello e dal colorito più scuro. Aveva una carnagione priva di
imperfezioni, escludendo le lentiggini che le circondavano gli occhi e
le davano uno sguardo ombroso: «La mia maschera da procione» diceva. La
cosa che più mi piaceva di lei - avevo raggiunto un'età in cui questi
dettagli assumevano un'importanza mal compresa ma innegabile - era il
suo sorriso. Sorrideva di rado ma in modo spettacolare. Era convinta di
avere i denti troppo sporgenti (si sbagliava) e aveva preso l'abitudine
di coprirsi la bocca quando rideva. Mi piaceva farla ridere, ma era il
suo sorriso che desideravo in segreto.

La settimana precedente, Jason aveva ricevuto dal padre un costoso
binocolo astronomico. Ci aveva giocherellato tutta la sera, puntandolo
sul poster turistico appeso sopra la TV e fingendo di spiare Cancún
dalla periferia di Washington. A un tratto però si alzò e disse:
«Dobbiamo andare a guardare il cielo.»

«No» ribatté subito Diane. «Si gela, fuori.»

«Ma è limpido. È la prima notte limpida della settimana. Fa solo un po'
freddo.»

«Stamattina c'era del ghiaccio sul prato.»

«Brina» la corresse lui.

«È mezzanotte passata.»

«È sabato sera.»

«Non possiamo uscire dal seminterrato.»

«Non possiamo disturbare la festa, nessuno ha parlato di non poter
andare fuori. Non ci vedrà nessuno, se hai paura di essere beccata.»

«Non ho paura di essere beccata.»

«E allora di \emph{cosa} hai paura?»

«Di sentirti parlare a vanvera mentre mi si congelano i piedi.»

Jason si girò verso di me. «E tu cosa dici, Tyler? Ti va di vedere un
po' di cielo?»

I gemelli mi chiedevano spesso di fare da arbitro alle loro discussioni,
con mio grande disagio. Era una proposta senza vincitori. Se mi fossi
schierato con Jason, mi sarei inimicato Diane e se mi fossi schierato
troppo spesso con Diane, sarebbe parso\ldots{} be', ovvio. Dissi: «Non
so, Jase, fa abbastanza freddo fuori\ldots»

Fu Diane a togliermi dall'impiccio. Mi mise una mano sulla spalla e
disse: «Non importa, un po' d'aria fresca sarà meglio che stare qui a
sentirlo lamentarsi.»

E così prendemmo le giacche dal corridoio del seminterrato e uscimmo
dalla porta sul retro.

La Grande Casa non era così imponente come poteva far credere il
soprannome che le avevamo dato, ma era più grande di una casa media di
quel quartiere benestante e si ergeva su un lotto di terra più vasto.
Un'ampia distesa ondeggiante di prati curati cedeva il passo, dietro di
sé, a un boschetto incolto di pini che delimitava un ruscello
leggermente inquinato. Per osservare le stelle, Jason scelse un punto a
metà strada fra la casa e il bosco.

Il mese di ottobre era stato gradevole fino al giorno prima, quando un
fronte freddo aveva tagliato le gambe all'estate indiana. Diane finse di
cingersi le costole e tremare ma era solo per punire Jason. L'aria
notturna era soltanto fresca, non era spiacevole. Il cielo era
cristallino e l'erba piuttosto asciutta, anche se forse la mattina dopo
sarebbe gelata di nuovo. La Luna non era visibile e non c'era traccia di
nuvole. La Grande Casa era illuminata come un vaporetto del Mississippi
e gettava il suo vivido bagliore giallo sul prato, ma sapevamo per
esperienza che nelle notti come quella, nascosti all'ombra di un albero,
saremmo stati completamente invisibili, come caduti in un buco nero.

Jason si sdraiò sulla schiena e puntò il binocolo al cielo stellato. Mi
sedetti con le gambe incrociate accanto a Diane e la vidi sfilare dalla
tasca della giacca una sigaretta, probabilmente rubata da sua madre.
(Carol Lawton, cardiologa ed ex fumatrice solo in teoria, nascondeva
pacchetti di sigarette nell'armadio, nella scrivania, in un cassetto
della cucina. Me l'aveva detto mia madre.) Se la mise in bocca, l'accese
con un accendino rosso trasparente - per un attimo la fiamma fu la cosa
più luminosa nello spazio intorno - ed esalò un pennacchio di fumo che
vorticò rapidamente nell'oscurità.

Mi sorprese a guardarla. «Vuoi fare un tiro?»

«Ha dodici anni» esclamò Jason. «Ha già i suoi problemi, ci manca solo
il cancro ai polmoni.»

«Certo» risposi. A quel punto era una questione d'onore.

Diane, divertita, mi passò la sigaretta. Aspirai con esitazione e
riuscii a non soffocare.

Se la riprese. «Non farti prendere la mano.»

«Tyler,» disse Jason, «sai qualcosa sulle stelle?»

Inspirai profondamente l'aria fredda e pulita. «Certo che sì.»

«Non mi riferisco alle cose che impari leggendo quei tascabili. Sai
farmi il \emph{nome} di qualche stella?»

Stavo arrossendo, ma speravo che il buio fosse sufficiente a non
farglielo notare. «Arturo» risposi. «Alfa Centauri. Sirio.
Polaris\ldots»

«E qual è il pianeta natale dei Klingon?»

«Non essere meschino» lo rimproverò Diane.

I gemelli erano davvero intelligenti per la loro età. Non che io fossi
stupido ma loro erano su un altro livello, e tutti ne eravamo
consapevoli. Loro frequentavano una scuola per bambini dotati, io andavo
alla scuola pubblica in autobus. Era una delle tante differenze evidenti
tra di noi. Loro vivevano nella Grande Casa, io stavo con mia madre in
un bungalow nella zona orientale della proprietà; i loro genitori si
dedicavano alla carriera, mia madre puliva loro la casa. In qualche modo
riconoscevamo questi divari senza farne un dramma.

«Okay,» disse Jason, «sai indicarmi Polaris?»

Polaris, la Stella del Nord. Avevo letto della schiavitù e della guerra
civile. C'era una canzone degli schiavi fuggitivi:

``\emph{Quando il Sole sorge, al primo verso della quaglia},

\emph{segui il Mestolino}.

\emph{Il vecchio ti aspetta per condurti alla libertà}

\emph{quando seguirai il Mestolino}.''

``\emph{Quando il Sole sorge}'' significava dopo il solstizio d'inverno.
La quaglia sverna a sud. Il mestolino era l'Orsa Maggiore, con
l'impugnatura puntata verso Polaris, a nord, la direzione della libertà:
trovai l'Orsa e speranzoso feci cenno con la mano in quella direzione.

«Vedi?» ribatté Diane a Jason, come se io avessi confermato un punto
essenziale di qualche discussione che non avevano voluto condividere con
me.

«Niente male» ammise Jason. «Sai cos'è una cometa?»

«Sì.»

«Vuoi vederne una?»

Annuii e mi stesi accanto a lui, gustando ancora e rimpiangendo il
sapore acre della sigaretta di Diane. Jason mi mostrò come puntellare i
gomiti sul terreno, poi mi fece avvicinare il binocolo agli occhi e
regolò il fuoco finché le stelle non divennero prima ovali indistinti e
poi puntini, molti più di quanti potessi vederne a occhio nudo.
Setacciai il cielo finché non trovai, o mi sembrò di trovare, il punto
che Jason mi aveva indicato: un minuscolo nodo di fosforescenza contro
il cielo spietatamente nero.

«Una cometa\ldots» cominciò a dire Jason.

«Lo so, una cometa è una specie di palla di neve polverosa che cade
verso il Sole.»

«Sì, possiamo dire così.» Era sprezzante. «Sai da dove vengono le
comete, Tyler? Vengono dal sistema solare esterno: da una specie di
alone ghiacciato attorno al Sole che si estende dall'orbita di Plutone
fino a metà strada rispetto alla stella più vicina, dove fa più freddo
di quanto tu possa immaginare.»

Annuii, un po' a disagio. Avevo letto abbastanza libri di fantascienza
da comprendere l'assoluta e indicibile grandezza del cielo notturno. A
volte mi piaceva pensarci, anche se farlo nel momento sbagliato della
notte, quando la casa era silenziosa, poteva mettermi addosso un po' di
paura.

«Diane?» chiese Jason. «Vuoi guardare?»

«Devo proprio?»

«No, certo che non devi. Puoi startene lì ad affumicarti i polmoni e
sbavare, se preferisci.»

«Intelligentone.» Spense la sigaretta nell'erba e tese la mano. Le
passai il binocolo.

«Stacci attenta a quello.» Jase stravedeva per il suo binocolo. Aveva
ancora l'odore dell'imballaggio di polistirolo e cellophane.

Lei regolò il fuoco e guardò verso l'alto. Rimase in silenzio per un
po'. Poi sbottò: «Sai cosa vedo quando uso questo coso per guardare le
stelle?»

«Cosa?»

«Le stesse stelle di sempre.»

«Usa l'immaginazione.» Jason sembrava davvero irritato.

«Se posso usare l'immaginazione perché dovrebbe servirmi un binocolo?»

«Pensa a quello che stai guardando, intendevo.»

«Ah» disse. Poi: «Oh. \emph{Oh}! Jason, vedo\ldots»

«Cosa?»

«Credo\ldots{} sì\ldots{} è Dio! E ha una lunga barba bianca! E regge un
cartello! E il cartello dice: J{ASON ROMPE}!»

«Molto divertente. Ridammelo se non sai usarlo.»

Lui tese la mano, lei lo ignorò, si mise a sedere e puntò il binocolo
verso le finestre della Grande Casa.

La festa continuava dal tardo pomeriggio. Mia madre mi aveva raccontato
che le feste dei Lawton erano ``costosi ritrovi per far scornare tra di
loro i pezzi da novanta delle imprese'' ma lei aveva un affinato senso
dell'iperbole per cui ciò che diceva andava ridimensionato. Quasi tutti
gli ospiti, aveva detto Jason, erano giovani promesse dell'industria
aerospaziale o membri di partito. Non la vecchia società di Washington
ma nuovi arrivati pieni di soldi con radici occidentali e agganci
nell'industria bellica. E.D. Lawton, il padre di Jason e Diane, ospitava
uno di questi eventi ogni tre o quattro mesi.

«È tutto come al solito» spiegò Diane da dietro gli ovali gemelli del
binocolo. «Al primo piano, ballano e bevono. In questo momento, bevono
più che ballare La cucina sembra chiusa, però. Penso che quelli del
catering si stiano preparando per andarsene. Tende dello studio tirate.
E.D. è nella biblioteca con un paio di manager. Che schifo! Uno di loro
fuma il sigaro.»

«Il tuo disgusto è poco convincente, signorina Marlboro» disse Jason.

Mentre lei continuava a elencare le finestre visibili, Jason si spostò e
mi venne accanto. «Tu le mostri l'universo,» bisbigliò, «e lei
preferisce spiare una cena.»

Non sapevo cosa rispondere. Come quasi tutto quello che Jason diceva,
suonò arguto e più intelligente di qualsiasi cosa avessi mai potuto
tirar fuori.

«La mia camera» spiegò Diane. «Vuota, grazie a Dio. La camera di Jason,
vuota a parte quella copia di Penthouse sotto al materasso\ldots»

«È un buon binocolo ma non fino a quel punto.»

«La camera di Carol e E.D., vuota; la stanza degli ospiti\ldots»

«Be'?»

Ma Diane non rispose. Rimase immobile con il binocolo davanti agli
occhi.

«Diane?» chiesi.

Restò in silenzio ancora per qualche secondo, poi rabbrividì, si voltò e
lanciò - \emph{scagliò} - il binocolo a Jason che protestò senza capire
che Diane aveva visto qualcosa di inquietante. Stavo per chiederle se
stesse bene\ldots{}

Quando le stelle scomparvero.

\separatorefiori

Non accadde un granché.

La gente lo dice spesso, la gente che l'ha visto succedere. Non accadde
un granché. Andò davvero così, e parlo da testimone: io stavo guardando
il cielo, mentre Diane e Jason bisticciavano. Non vidi nulla se non uno
strano bagliore, che mi lasciò negli occhi un'immagine residua delle
stelle di un verde freddo e fluorescente. Battei le palpebre.

Jason chiese: «Cos'era quello? Un fulmine?»

Diane invece rimase zitta.

«Jason» mormorai continuando a battere le palpebre.

«Cosa? Diane, giuro su Dio, se hai spaccato una lente\ldots»

«Sta' zitto» strillò lei.

E io: «Smettetela. \emph{Guardate}. Cos'è successo alle stelle?»

Entrambi alzarono la testa verso il cielo.

\separatorefiori

Di noi tre, solo Diane era disposta a credere che le stelle fossero
effettivamente ``morte'', che si fossero spente come candele al vento.
Era impossibile, sosteneva Jason: la luce di quelle stelle aveva
percorso cinquanta anni luce, o cento, o cento milioni, a seconda della
sorgente; sicuramente non avevano smesso di brillare tutte insieme in
una sequenza enormemente elaborata e ideata per farla sembrare
simultanea ai terrestri. Comunque, feci notare, anche il Sole era una
stella, e stava \emph{ancora} brillando, almeno dall'altra parte del
pianeta, o no? Certo che sì. Altrimenti, disse Jason, entro l'indomani
mattina saremmo tutti morti congelati.

Di conseguenza, a rigor di logica, le stelle brillavano ancora ma noi
non riuscivamo a vederle. Non erano sparite, ma oscurate: eclissate. Sì,
il cielo all'improvviso era diventato un vuoto nero come l'ebano, però
era un mistero, non una catastrofe.

Ma un altro aspetto del commento di Jason mi si era ficcato in testa. E
se il Sole fosse davvero scomparso? Mi immaginai la neve cadere leggera
nell'oscurità perenne, e poi supposi che l'aria stessa si sarebbe
congelata in un tipo diverso di neve fino a seppellire tutta la civiltà
umana sotto l'oggetto del suo stesso respiro. Ecco perché era meglio,
senza dubbio meglio, ipotizzare che le stelle fossero state
``eclissate''. Ma da cosa?

«Be', ovviamente da qualcosa di grande, qualcosa di veloce. Tu l'hai
visto, Tyler. È successo tutto a un tratto o hai notato qualcosa che si
muoveva nel cielo?»

Gli dissi che era come se le stelle avessero brillato e poi si fossero
spente tutte insieme.

«Fanculo quelle stupide stelle» sibilò Diane. (Rimasi scioccato:
\emph{fanculo} non era una parola che usava abitualmente, cosa che io e
Jase invece facevamo liberamente ora che la nostra età aveva raggiunto
da un po' le due cifre. Quell'estate erano cambiate molte cose.)

Jason percepì dell'ansia nella sua voce. «Non credo ci sia niente da
temere» disse, sebbene lui stesso fosse visibilmente turbato.

Diane si accigliò. «Sento freddo.»

E così decidemmo di tornare alla Grande Casa per vedere se la notizia
era arrivata alla CNN o CNBC. Mentre attraversavamo il prato il cielo
era inquietante, completamente nero, senza peso ma opprimente, più scuro
di qualsiasi cielo avessi mai visto.

\separatorefiori

«Dobbiamo dirlo a E.D.» fece Jason.

«Glielo dici tu» rispose Diane.

Jase e Diane chiamavano i loro genitori per nome, perché Carol Lawton
immaginava che la sua fosse una famiglia progressista. La realtà era ben
più complessa. Carol era indulgente ma non si lasciava coinvolgere molto
dalla vita dei gemelli, mentre E.D. agiva sistematicamente con lo scopo
di tirare su un erede. Quell'erede, naturalmente, era Jason. Jason
venerava suo padre. Diane lo temeva. E io non ero così sprovveduto da
farmi vedere nella zona degli adulti alla fine di un ricevimento dei
Lawton dove tutti ormai erano ubriachi; e così io e Diane indugiammo
nella zona demilitarizzata dietro una porta, mentre Jason trovò suo
padre in una stanza adiacente. Non riuscimmo a sentire nei dettagli la
loro conversazione ma il tono di voce di E.D. era inequivocabile:
risentito, insofferente e irritato. Jason ritornò nel seminterrato con
la faccia rossa e quasi in lacrime, io mi scusai e mi diressi verso la
porta sul retro.

Diane mi raggiunse nel corridoio. Mi appoggiò la mano sul polso come se
volesse ancorarci l'uno all'altra. «Tyler,» disse, «si alzerà, vero? Il
Sole, intendo, domattina. So che è una domanda stupida, ma il Sole
\emph{sorgerà}, giusto?»

Sembrava davvero disperata. Cominciai a dire qualcosa di superficiale -
«se non sorge, moriremo tutti» - ma la sua ansia fece dubitare anche me.
Che cosa avevamo visto di preciso e che significava? Jason,
evidentemente, non aveva saputo convincere suo padre del fatto che
qualcosa d'importante era accaduto nel cielo notturno, perciò forse ci
stavamo sgomentando per niente. E se invece il mondo era davvero alla
fine e solo noi tre lo sapevamo?

«Staremo bene» la rassicurai.

Mi scrutò attraverso la palizzata di capelli lisci. «Tu ci credi?»

Tentai di sorridere. «Al novanta per cento.»

«Ma rimarrai sveglio fino a domattina, vero?»

«Forse, probabilmente.» Sapevo già di non voler dormire.

Fece un gesto con il pollice e il mignolo: «Posso chiamarti più tardi?»

«Certo.»

«Probabilmente non riuscirò a dormire. Lo so che sembra sciocco ma, in
caso mi addormentassi, mi chiami appena sorge il Sole?»

Dissi che lo avrei fatto.

«Promesso?»

«Promesso.» Ero elettrizzato che me lo avesse chiesto.

\separatorefiori

La casa in cui vivevo con mia madre era un bungalow in legno ben tenuto,
situato nella parte orientale della proprietà dei Lawton. Un piccolo
giardino di rose recintato da una staccionata di pino stava addossato ai
gradini d'ingresso - ma le rose, rifiorite quell'autunno, erano
appassite con l'ultima ventata di aria fredda. In quella notte senza
Luna, nuvole e stelle, la luce del portico scintillava come un faro.

Entrai senza fare rumore. Mia madre si era ritirata in camera sua già da
un po'. Il piccolo salotto era in ordine, salvo per un unico bicchierino
vuoto sulla credenza: nei giorni feriali era astemia, mentre nei fine
settimana si concedeva un po' di whisky. Diceva di avere solo due vizi,
e un bicchiere il sabato sera era uno di questi. (Una volta, quando le
chiesi quale fosse l'altro, mi guardò a lungo e rispose: «Tuo padre.»
Non insistetti sull'argomento.)

Mi sdraiai sul divano vuoto a leggere fino a quando non mi chiamò Diane,
meno di un'ora dopo. La prima cosa che disse fu: «Hai acceso la TV?»

«Dovrei?»

«Non scomodarti, non trasmettono niente.»

«Be', sai com'è, sono le \emph{due} del mattino.»

«No, intendevo proprio niente di niente. Ci sono solo le televendite sui
canali locali, per il resto nulla. Che significa, Tyler?»

Significava che ogni satellite in orbita era svanito insieme alle
stelle. Le telecomunicazioni, il meteo, i satelliti militari, il sistema
GPS: erano stati tutti interrotti in un batter d'occhio. Ma non sapevo
nulla di tutto ciò e di certo non sarei stato in grado di spiegarlo a
Diane. «Potrebbe voler dire qualsiasi cosa.»

«Fa un po' paura.»

«Probabilmente niente di cui preoccuparsi.»

«Lo spero. Sono contenta che tu sia ancora sveglio.»

Richiamò un'ora dopo per darmi altre notizie. Anche la connessione
internet non funzionava, mi spiegò. E la TV locale aveva cominciato a
riferire di voli mattutini annullati al Reagan e negli aeroporti
regionali, avvertendo le persone di chiamare prima.

«Ma tutta la notte hanno volato dei jet.» Ne avevo visto le luci dalla
finestra della mia camera: false stelle, veloci. «Militari, immagino.
Potrebbe essere roba di terrorismo.»

«Jason è in camera sua con una radio, si sta sintonizzando sulle
stazioni di Boston e New York. Dice che parlano di attività militari e
chiusura degli aeroporti ma non di terrorismo - e neanche di stelle.»

«Qualcuno deve averlo notato.»

«Anche se fosse, non ne fanno cenno. Forse hanno ricevuto \emph{ordini}
di non farlo. Non hanno detto nulla neanche dell'alba.»

«E perché avrebbero dovuto? Il Sole sorgerà fra quanto? Un'ora? Vuol
dire che sull'oceano sta già salendo, al largo della costa atlantica. Le
navi in mare devono averlo visto. Lo \emph{vedremo}, fra non molto.»

«Lo spero.» Sembrava spaventata e imbarazzata allo stesso tempo. «Spero
che tu abbia ragione.»

«Vedrai.»

«Mi piace la tua voce, Tyler. Te l'ho mai detto? Hai una voce molto
rassicurante.»

Anche se quelle che avevo raccontato erano solo stronzate.

Ma quel complimento mi colpì molto più di quanto volessi farle
intendere. Ci pensai dopo che ebbe riattaccato. Continuai a ripetermelo
in mente per il gusto di sentire la sensazione di calore che provocava.
E mi chiedevo cosa volesse dire. Diane aveva un anno più di me ma era
tre volte più sofisticata - e allora perché all'improvviso mi sentivo
così protettivo nei suoi confronti, perché desideravo averla vicino
tanto da toccarle il viso e prometterle che tutto sarebbe andato bene?
Era un enigma impellente e inquietante, quasi quanto ciò che era
accaduto al cielo.

\separatorefiori

Richiamò alle cinque meno dieci quando, mio malgrado, mi ero quasi
appisolato completamente vestito. Cercai a tastoni il telefono nella
tasca della camicia. «Pronto?»

«Sono io. È ancora buio, Tyler.»

Gettai un'occhiata alla finestra. Sì. Buio. Poi all'orologio sul
comodino. «Non è ancora l'alba, Diane.»

«Dormivi?»

«No.»

«Sì che dormivi. Beato te. È ancora buio. Fa anche freddo. Ho guardato
il termometro fuori dalla finestra della cucina. Due gradi. È normale
che faccia tanto freddo?»

«Pure ieri mattina faceva così freddo. C'è qualcun altro sveglio, da
te?»

«Jason si è chiuso in camera sua con la radio. I miei, ehm, genitori,
ehm, credo che stiano smaltendo la festa. Tua mamma è sveglia?»

«Non a quest'ora, non nel fine settimana.» Lanciai un'altra occhiata
nervosa alla finestra. Di certo, a quell'ora, in cielo avrebbe dovuto
esserci un po' di luce. Sarebbe stato rassicurante anche un accenno di
alba.

«Non l'hai svegliata?»

«E cosa dovrebbe fare, Diane? Far tornare le stelle?»

«Credo di no.» Fece una pausa. «Tyler» mormorò.

«Sono ancora qui.»

«Qual è la prima cosa che ricordi?»

«Intendi\ldots{} di oggi?»

«No. La prima cosa che ricordi in assoluto della tua vita. So che è una
domanda stupida, ma credo che starò meglio se per cinque o dieci minuti
parliamo di qualcosa che non sia il cielo.»

«La prima cosa?» Ci pensai un po' su. «Direi quando ero a Los Angeles,
prima che ci trasferissimo a est.» Quando mio padre era ancora vivo e
lavorava ancora per E.D. Lawton nell'azienda appena avviata a
Sacramento. «Avevamo un appartamento con delle grosse tende bianche
nella camera da letto. La prima cosa che ricordo davvero è quando
guardavo quelle tende scosse dal vento. Era una giornata soleggiata, la
finestra era aperta e soffiava la brezza.» Inaspettatamente quel ricordo
mi commosse, come l'ultima immagine di una battigia che spariva. «E tu,
invece?»

La prima cosa che Diane rammentava era un momento che, come me, aveva
trascorso a Sacramento, solo che era molto diverso dal mio. E.D. aveva
portato i figli a fare un giro dello stabilimento, assegnando già da
allora a Jason il ruolo di erede favorito. Diane era rimasta affascinata
dagli enormi pali perforati sul pavimento della fabbrica, dalle bobine
di tessuto di alluminio ultrasottile grandi come case, dal rumore
costante. Era tutto tanto enorme che si aspettava quasi di trovare un
gigante delle fiabe incatenato alle mura, prigioniero di suo padre.

Non era un bel ricordo. Disse che si era sentita trascurata, quasi
persa, abbandonata dentro un enorme e terrificante macchinario da
costruzione.

Ne parlammo per un po'. Poi Diane disse: «Controlla il cielo.»

Guardai verso la finestra. La luce che si riversava sull'orizzonte
occidentale era sufficiente a trasformare l'oscurità in un blu
inchiostro.

Non volli confessare il sollievo che provai.

«Credo che avessi ragione» disse lei, a un tratto briosa. «Alla fine, il
Sole sta sorgendo.»

Naturalmente non era il vero Sole. Era un astro impostore, una
falsificazione ingegnosa, ma noi ancora non lo sapevamo.